% --- Archon physics formalization source begin ---
% archon:physics
% archon:covers PhyXMiniProblems/problem_phyx_mini_0098.lean
% archon:source-report reports/phyx_mini/problem_phyx_mini_0098.source.json

\chapter{Physics problem phyx\_mini\_0098}
\label{ch:PhyXMiniProblems_problem_phyx_mini_0098}

\paragraph{Problem source.}
\#\# Physical scenario

Optical fibers are constructed with a cylindrical core surrounded by a sheath of cladding material. Common materials used are pure silica (\$n\_2 = 1.450\$) for the cladding and silica doped with germanium (\$n\_1 = 1.465\$) for the core.

\#\# Question

What is the value of \$	heta\_i\$ for \$n\_1 = 1.465\$ and \$n\_2 = 1.450\$?

\#\# Answer choices

- A: A: 12.1°

- B: B: 13.1°

- C: C: 11.9°

- D: D: 13.2°

\#\# Figure caption (auxiliary; use the image as primary evidence)

The image appears to represent optical fiber communication. It depicts a cylindrical fiber optic cable with two layers: an inner core and an outer cladding. 

- The inner core is labeled with the refractive index \textbackslash{}( n\_1 \textbackslash{}).
- The outer cladding is labeled with the refractive index \textbackslash{}( n\_2 \textbackslash{}).
- A purple arrow represents a light ray entering the fiber at an angle \textbackslash{}(\textbackslash{}theta\_i\textbackslash{}) relative to the normal line (dashed).
- The light ray undergoes total internal reflection, bouncing along the core-cladding interface.

The diagram illustrates how light travels through the fiber due to differences in refractive indices, with \textbackslash{}( n\_1 > n\_2 \textbackslash{}), ensuring containment of light within the core.

\#\# Dataset metadata

Geometrical Optics; Physical Model Grounding Reasoning, Spatial Relation Reasoning

\paragraph{Recorded answer/context.}
A: A: 12.1°

\paragraph{Figure/image path.}
/root/proposal\_for\_physic/science-mango/phyx\_mini\_run/phyx\_data/test\_image/98.png

\paragraph{Formalization target.}
create a compiling Lean file with sorry bodies at `PhyXMiniProblems/problem\_phyx\_mini\_0098.lean`. The Lean declarations must preserve the physical quantities, dimensions or dimensional roles, figure labels, governing-law hypotheses, and final relation expressed by this problem.
Use Mathlib/Physlib names found through LeanExplore where available. If a physics API is missing, introduce faithful local abstractions rather than scalar placeholder aliases.

\begin{theorem}[Physics formalization target]
\label{thm:physics:phyx_mini_0098:target}
\lean{PhyXMiniProblems.ProblemPhyXMini0098.theta_i_matches_choice_A}
\uses{def:physics:phyx-mini-0098:phyxminiproblems-problemphyxmini0098-stepindexfibersetup, def:physics:phyx-mini-0098:phyxminiproblems-problemphyxmini0098-matchesopticalfiberfigure, def:physics:phyx-mini-0098:phyxminiproblems-problemphyxmini0098-satisfieslimitingguidancelaws, def:physics:phyx-mini-0098:phyxminiproblems-problemphyxmini0098-matchesnearesttenthdegree}
The assigned autoformalize agent should translate this physics problem into Lean declarations in the covered file, with theorem and lemma proof bodies written as `by sorry`.
\end{theorem}
\begin{proof}
This is an autoformalization task, not a proof task. Produce statements that can later be proved by the physics prover without weakening the physical model.
\end{proof}

% --- Archon declaration topology begin ---
\section{Declaration topology}
\begin{definition}[Dimensionless Optical Index]
  \label{def:physics:phyx-mini-0098:phyxminiproblems-problemphyxmini0098-dimensionlessopticalindex}
  \lean{PhyXMiniProblems.ProblemPhyXMini0098.DimensionlessOpticalIndex}
  A unit-independent physical scalar with the identity (dimensionless) dimension
\end{definition}
\begin{proof}
  This is the declared model definition of Dimensionless Optical Index.
\end{proof}
\begin{definition}[refractive Index Readout]
  \label{def:physics:phyx-mini-0098:phyxminiproblems-problemphyxmini0098-refractiveindexreadout}
  \lean{PhyXMiniProblems.ProblemPhyXMini0098.refractiveIndexReadout}
  \uses{def:physics:phyx-mini-0098:phyxminiproblems-problemphyxmini0098-dimensionlessopticalindex}
  The real SI readout of a dimensionless optical refractive index
\end{definition}
\begin{proof}
  This is the declared model definition of refractive Index Readout.
\end{proof}
\begin{definition}[Optical Fiber Region]
  \label{def:physics:phyx-mini-0098:phyxminiproblems-problemphyxmini0098-opticalfiberregion}
  \lean{PhyXMiniProblems.ProblemPhyXMini0098.OpticalFiberRegion}
  The three homogeneous optical regions traversed by the depicted ray
\end{definition}
\begin{proof}
  This is the declared model definition of Optical Fiber Region.
\end{proof}
\begin{definition}[degrees]
  \label{def:physics:phyx-mini-0098:phyxminiproblems-problemphyxmini0098-degrees}
  \lean{PhyXMiniProblems.ProblemPhyXMini0098.degrees}
  Convert a numerical degree readout into a physical angle
\end{definition}
\begin{proof}
  This is the declared model definition of degrees.
\end{proof}
\begin{definition}[degree Readout]
  \label{def:physics:phyx-mini-0098:phyxminiproblems-problemphyxmini0098-degreereadout}
  \lean{PhyXMiniProblems.ProblemPhyXMini0098.degreeReadout}
  Express the acute representative of a physical angle in degrees
\end{definition}
\begin{proof}
  This is the declared model definition of degree Readout.
\end{proof}
\begin{definition}[Step Index Fiber Setup]
  \label{def:physics:phyx-mini-0098:phyxminiproblems-problemphyxmini0098-stepindexfibersetup}
  \lean{PhyXMiniProblems.ProblemPhyXMini0098.StepIndexFiberSetup}
  \uses{def:physics:phyx-mini-0098:phyxminiproblems-problemphyxmini0098-dimensionlessopticalindex, def:physics:phyx-mini-0098:phyxminiproblems-problemphyxmini0098-opticalfiberregion}
  The dimensionless optical data and angles of the limiting guided ray. thetaI is the figure label θᵢ, measured from the dashed fiber axis (and hence from the entrance-face normal). Inside the core, the same ray makes coreRayAngleToAxis with the axis and coreCladdingIncidenceAngle with the normal to the cylindrical core--cladding interface
\end{definition}
\begin{proof}
  This is the declared model definition of Step Index Fiber Setup.
\end{proof}
\begin{definition}[Matches Optical Fiber Figure]
  \label{def:physics:phyx-mini-0098:phyxminiproblems-problemphyxmini0098-matchesopticalfiberfigure}
  \lean{PhyXMiniProblems.ProblemPhyXMini0098.MatchesOpticalFiberFigure}
  \uses{def:physics:phyx-mini-0098:phyxminiproblems-problemphyxmini0098-refractiveindexreadout, def:physics:phyx-mini-0098:phyxminiproblems-problemphyxmini0098-stepindexfibersetup}
  Numerical material readouts from the problem, together with the usual idealization that the unlabelled exterior region is air of index one
\end{definition}
\begin{proof}
  This is the declared model definition of Matches Optical Fiber Figure.
\end{proof}
\begin{definition}[Is Acute Optical Angle]
  \label{def:physics:phyx-mini-0098:phyxminiproblems-problemphyxmini0098-isacuteopticalangle}
  \lean{PhyXMiniProblems.ProblemPhyXMini0098.IsAcuteOpticalAngle}
  Select the acute representative appropriate for an angle in the diagram
\end{definition}
\begin{proof}
  This is the declared model definition of Is Acute Optical Angle.
\end{proof}
\begin{definition}[Satisfies Limiting Guidance Laws]
  \label{def:physics:phyx-mini-0098:phyxminiproblems-problemphyxmini0098-satisfieslimitingguidancelaws}
  \lean{PhyXMiniProblems.ProblemPhyXMini0098.SatisfiesLimitingGuidanceLaws}
  \uses{def:physics:phyx-mini-0098:phyxminiproblems-problemphyxmini0098-refractiveindexreadout, def:physics:phyx-mini-0098:phyxminiproblems-problemphyxmini0098-opticalfiberregion, def:physics:phyx-mini-0098:phyxminiproblems-problemphyxmini0098-degrees, def:physics:phyx-mini-0098:phyxminiproblems-problemphyxmini0098-stepindexfibersetup, def:physics:phyx-mini-0098:phyxminiproblems-problemphyxmini0098-isacuteopticalangle}
  The geometrical-optics laws for the limiting ray accepted by the fiber. At the entrance face the ray obeys Snell's law. The normal to the cylindrical wall is perpendicular to the fiber axis, so the two internal angles are complementary. At the limiting acceptance angle, transmission into the cladding would be tangential; the final field is precisely Snell's law for that critical condition. None of these laws contains the requested one-decimal answer
\end{definition}
\begin{proof}
  This is the declared model definition of Satisfies Limiting Guidance Laws.
\end{proof}
\begin{definition}[Answer Choice]
  \label{def:physics:phyx-mini-0098:phyxminiproblems-problemphyxmini0098-answerchoice}
  \lean{PhyXMiniProblems.ProblemPhyXMini0098.AnswerChoice}
  Labels of the four displayed multiple-choice answers
\end{definition}
\begin{proof}
  This is the declared model definition of Answer Choice.
\end{proof}
\begin{definition}[angle Degrees]
  \label{def:physics:phyx-mini-0098:phyxminiproblems-problemphyxmini0098-answerchoice-angledegrees}
  \lean{PhyXMiniProblems.ProblemPhyXMini0098.AnswerChoice.angleDegrees}
  \uses{def:physics:phyx-mini-0098:phyxminiproblems-problemphyxmini0098-answerchoice}
  The degree value printed beside each answer choice
\end{definition}
\begin{proof}
  This is the declared model definition of angle Degrees.
\end{proof}
\begin{definition}[Matches Nearest Tenth Degree]
  \label{def:physics:phyx-mini-0098:phyxminiproblems-problemphyxmini0098-matchesnearesttenthdegree}
  \lean{PhyXMiniProblems.ProblemPhyXMini0098.MatchesNearestTenthDegree}
  \uses{def:physics:phyx-mini-0098:phyxminiproblems-problemphyxmini0098-degreereadout, def:physics:phyx-mini-0098:phyxminiproblems-problemphyxmini0098-answerchoice}
  A displayed one-decimal choice matches a physical angle when its degree readout differs by at most half of one tenth of a degree
\end{definition}
\begin{proof}
  This is the declared model definition of Matches Nearest Tenth Degree.
\end{proof}
\begin{lemma}[acceptance angle numerical aperture relation]
  \label{lem:physics:phyx-mini-0098:phyxminiproblems-problemphyxmini0098-acceptance-angle-numerical-aperture-relation}
  \lean{PhyXMiniProblems.ProblemPhyXMini0098.acceptance_angle_numerical_aperture_relation}
  \uses{def:physics:phyx-mini-0098:phyxminiproblems-problemphyxmini0098-refractiveindexreadout, def:physics:phyx-mini-0098:phyxminiproblems-problemphyxmini0098-stepindexfibersetup, def:physics:phyx-mini-0098:phyxminiproblems-problemphyxmini0098-satisfieslimitingguidancelaws}
  The entrance and critical-interface laws imply the standard numerical-aperture relation. This intermediate result contains no answer-choice value
\end{lemma}
\begin{proof}
  The conclusion follows from the governing relations and figure readouts named in the statement of acceptance angle numerical aperture relation.
\end{proof}
% --- Archon declaration topology end ---
% --- Archon physics formalization source end ---
